
En el presente avance (PDF~5), la aplicaci\'on m\'ovil \nombreApp{} ha
alcanzado un nivel de madurez significativo en t\'erminos de persistencia
de datos y operaciones CRUD. Las principales conclusiones del avance son
las siguientes:

\begin{enumerate}
  \item \textbf{An\'alisis completo de datos cr\'iticos.} Se identificaron
    seis conjuntos de datos que requieren persistencia (sesi\'on,
    favoritos, escenarios, eventos, perfil e im\'agenes) y se clasificaron
    por criticidad, ciclo de vida y nivel de sensibilidad.

  \item \textbf{Mecanismos de persistencia seleccionados.} Se adopt\'o una
    arquitectura de tres niveles: AsyncStorage para la sesi\'on de usuario,
    Supabase (PostgreSQL) para los datos de dominio, y React Query para
    el cach\'e en memoria. Esta combinaci\'on cubre las necesidades de
    persistencia local, remota y de rendimiento.

  \item \textbf{Estructura de base de datos dise~nada.} Se documentaron
    las siete tablas principales del sistema (profiles, scenarios,
    scenario\_images, sports, scenario\_sports, events, favorites), sus
    relaciones y las pol\'iticas RLS que garantizan la seguridad de los
    datos a nivel de fila.

  \item \textbf{CRUD completo implementado.} Se implementaron las 13
    operaciones CRUD para las tres entidades principales (escenarios,
    eventos y favoritos), incluyendo operaciones de soft delete, hard
    delete y upsert. Todas las operaciones se sincronizan mediante
    invalidaci\'on de cach\'e de React Query.

  \item \textbf{Persistencia tras el cierre verificada.} Se verific\'o
    que la sesi\'on del usuario se conserva mediante AsyncStorage, que
    los favoritos persisten en Supabase y que los escenarios y eventos
    se recuperan autom\'aticamente al reabrir la aplicaci\'on.

  \item \textbf{Optimizaci\'on del rendimiento.} Se crearon \'indices
    en las columnas m\'as consultadas de la base de datos, reduciendo
    el tiempo de respuesta de las consultas principales. El cach\'e de
    React Query con staleTime de 5 minutos evita consultas redundantes.

  \item \textbf{Seguridad reforzada.} Las pol\'iticas RLS garantizan
    que cada usuario solo pueda acceder a sus propios datos. La
    eliminaci\'on de escenarios y eventos est\'a restringida al rol
    de administrador, y la subida de im\'agenes tambi\'en requiere
    permisos de administrador.

  \item \textbf{Gesti\'on de escenarios completa.} La nueva pantalla
    de gesti\'on (\textit{manage.tsx}) permite a los administradores y
    gestores crear, editar y eliminar escenarios deportivos, incluyendo
    la subida de im\'agenes y la gesti\'on de eventos asociados.
\end{enumerate}

La aplicaci\'on ahora cuenta con una capa de persistencia robusta que
garantiza la conservaci\'on de los datos cr\'iticos del usuario y del
sistema. Los datos se almacenan de forma segura tanto en el dispositivo
(AsyncStorage) como en el servidor (Supabase), y se sincronizan de
forma autom\'atica a trav\'es de React Query. El CRUD completo permite
la gesti\'on integral de escenarios, eventos y favoritos, estableciendo
las bases para funcionalidades futuras como la gesti\'on de usuarios y
los reportes estad\'isticos.
